Composing (2.1.2) and (2.1.3), we obtain the homomorphisms (2.1.1).

The algebra structure defined by the homomorphisms (2.1.1) is associative (cf. [1] XI 2). Furthermore, the @@M0@@-module @@M1@@ has a standard @@M2@@-module structure, with @@M3@@ acting in the evident way, and @@M4@@ by functoriality; likewise, the @@M5@@ carry two @@M6@@-module structures. The homomorphisms (2.1.2) and (2.1.3) are clearly compatible with these two structures: for (2.1.2), by functoriality (working in the category of @@M7@@-modules), and for (2.1.3), by the classical result on morphisms of an acyclic complex. We have therefore endowed @@M8@@ with an associative @@M9@@-algebra structure.

\subsection*{2.2.}
Suppose that @@M10@@ admits a resolution by projective @@M11@@-modules @@M12@@ such that @@M13@@ is a graded @@M14@@-algebra, with the differential satisfying:
@@M15@@
\SGAAuthorFootText{1}{One can show that such a resolution always exists.}
Then the ``product'' maps define commutative diagrams for @@M16@@:
\sgaStableTarget{sga6.e07.eq.2-2-1}@@M17@@
and consequently a morphism of complexes @@M18@@ extending the augmentation homomorphism: @@M19@@. This shows that the algebra structure of @@M20@@ can be computed from that of @@M21@@. In particular, if @@M22@@ is an alternating @@M23@@-algebra (i.e. such that for every
